% ============================================================
% CAPÍTULO: ESTIMACIÓN EMPÍRICA
% Tesis doctoral: Eficiencia hospitalaria, estructura
% organizativa y teoría de agencia
% José María Elena Izquierdo — Universidad de Salamanca
% Marzo 2026
% ============================================================

\documentclass[12pt, a4paper]{article}

% ── Paquetes ─────────────────────────────────────────────────
\usepackage[utf8]{inputenc}
\usepackage[T1]{fontenc}
\usepackage[spanish]{babel}
\usepackage{amsmath, amssymb, amsthm}
\usepackage{booktabs}
\usepackage{longtable}
\usepackage{array}
\usepackage{multirow}
\usepackage{graphicx}
\usepackage{geometry}
\usepackage{setspace}
\usepackage{natbib}
\usepackage{hyperref}
\usepackage{caption}
\usepackage{subcaption}
\usepackage{threeparttable}
\usepackage{rotating}
\usepackage{pdflscape}
\usepackage{xcolor}
\usepackage{lmodern}
\usepackage{microtype}
\usepackage{footnote}
\usepackage{chngcntr}

\geometry{
  top=2.5cm, bottom=2.5cm,
  left=3cm, right=2.5cm
}
\onehalfspacing

\hypersetup{
  colorlinks=true,
  linkcolor=blue!60!black,
  citecolor=blue!60!black,
  urlcolor=blue!60!black
}

% ── Comandos propios ─────────────────────────────────────────
\newcommand{\E}{\mathbb{E}}
\newcommand{\N}{\mathcal{N}}
\newcommand{\bsym}[1]{\boldsymbol{#1}}
\newcommand{\sig}[1]{\ensuremath{^{#1}}}
\newcommand{\note}[1]{\textit{\small #1}}

\counterwithin{table}{section}
\counterwithin{figure}{section}

% ── Información del documento ────────────────────────────────
\title{
  \large Tesis Doctoral\\[0.3em]
  \LARGE\textbf{Eficiencia hospitalaria, estructura organizativa\\
  y teoría de agencia}\\[0.5em]
  \large\textit{Capítulo de Estimación Empírica}\\[0.3em]
  \normalsize Documento de trabajo para revisión del
  director de tesis
}
\author{José María Elena Izquierdo\\
  \small Departamento de Economía Aplicada\\
  \small Universidad de Salamanca}
\date{Marzo 2026}

% ═══════════════════════════════════════════════════════════
\begin{document}
% ═══════════════════════════════════════════════════════════

\maketitle
\thispagestyle{empty}
\newpage

\tableofcontents
\newpage

% ═══════════════════════════════════════════════════════════
\section{Introducción y motivación empírica}
% ═══════════════════════════════════════════════════════════

Este capítulo desarrolla la estrategia empírica para
contrastar las predicciones del modelo de agencia
multitarea de \citet{izquierdo2019}. El modelo describe
la relación entre el pagador público (Sistema Nacional de
Salud, SNS), el hospital y el médico, y genera
predicciones contrastables sobre cómo la estructura
organizativa y el mecanismo de pago afectan de forma
\textit{asimétrica} a dos dimensiones del output
hospitalario: la cantidad de atención prestada y su
intensidad diagnóstico-terapéutica.

La contribución empírica central es identificar y
contrastar esta asimetría utilizando un sistema de
fronteras estocásticas (SFA) con datos de panel del
sistema hospitalario español para el período 2010--2023.
La evidencia producida es, hasta donde conocemos, la
primera contrastación empírica directa de un modelo de
agencia multitarea en el sector hospitalario español
mediante análisis de frontera estocástica.

El capítulo se organiza del siguiente modo. La sección~2
presenta el modelo teórico y sus predicciones en forma
contrastable. La sección~3 describe la estrategia
empírica y los cuatro diseños de estimación. La
sección~4 detalla los datos y las variables. La
sección~5 presenta los resultados principales. La
sección~6 reporta los análisis de robustez y los tests
de especificación. La sección~7 discute las implicaciones
de los resultados y sus limitaciones.

% ═══════════════════════════════════════════════════════════
\section{El modelo teórico y sus predicciones contrastables}
% ═══════════════════════════════════════════════════════════

\subsection{Estructura del modelo}

El modelo de \citet{izquierdo2019} es un juego de
principal-agente en tres niveles con dos outputs
producidos asimétricamente. El pagador (SNS) contrata
con el hospital, que a su vez contrata con el médico.
Los outputs son:

\begin{itemize}
  \item \textbf{$q$ (cantidad)}: número de altas
    ponderadas por complejidad. Resulta del esfuerzo
    conjunto del hospital ($a$) y del médico ($t_1$):
    $q = q(a, t_1)$.
  \item \textbf{$i$ (intensidad)}: nivel de recursos
    diagnóstico-terapéuticos aplicados por episodio.
    Depende exclusivamente del esfuerzo médico ($t_2$):
    $i = i(t_2)$.
\end{itemize}

Esta asimetría en la producción es el núcleo estructural
del modelo. El parámetro $\psi_{12}$ mide la
complementariedad o sustituibilidad entre las dos tareas
del médico:
\begin{equation}
  \psi_{12} = \frac{\partial^2 C}{\partial t_1 \partial t_2}
  \label{eq:psi12}
\end{equation}
donde $C(\cdot)$ es la función de coste del esfuerzo
médico. Si $\psi_{12} > 0$ (tareas sustitutas), el médico
que dedica más tiempo a ver pacientes reduce la
intensidad por episodio --- caso típico de servicios
médicos no quirúrgicos. Si $\psi_{12} < 0$ (tareas
complementarias), el acto de intervenir genera per se
mayor intensidad --- caso típico de servicios quirúrgicos.

\subsection{Estructuras organizativas}

El modelo distingue dos configuraciones institucionales:

\begin{description}
  \item[Centralizada ($D^{desc} = 0$):] Hospitales SNS
    con personal estatutario. El pagador contrata
    simultáneamente con el hospital y el médico, manteniendo
    control directo sobre ambos niveles y pagando mediante
    presupuesto global retrospectivo.
  \item[Descentralizada ($D^{desc} = 1$):] Concesiones,
    fundaciones sanitarias y hospitales privados
    concertados. El pagador contrata únicamente con el
    hospital, que diseña autónomamente el contrato con el
    médico y asume la responsabilidad sobre los incentivos
    internos. El pago al hospital es, en general, de
    carácter prospectivo por actividad.
\end{description}

\subsection{Predicciones contrastables}

El equilibrio del modelo genera cuatro predicciones
sobre los determinantes de la ineficiencia técnica en
cada dimensión del output (véase Tabla~\ref{tab:predicciones}):

\begin{table}[h!]
\centering
\caption{Predicciones contrastables del modelo teórico}
\label{tab:predicciones}
\begin{threeparttable}
\begin{tabular}{clcc}
\toprule
\textbf{Pred.} & \textbf{Descripción} &
\textbf{Parámetro} & \textbf{Signo esperado} \\
\midrule
P1 & Descentralización $\downarrow$ ineficiencia en $q$
   & $\alpha_1$ en $\mu^q_{it}$ & $\alpha_1 < 0$ \\
P2 & Descentralización $\uparrow$ ineficiencia en $i$
   & $\delta_1$ en $\mu^i_{it}$ & $\delta_1 > 0$ \\
P3 & El tipo de pago difiere entre dimensiones
   & $\alpha_2 \neq \delta_2$ & Signos opuestos \\
P4 & $\psi_{12}$ modera asimétricamente el diferencial
   & $\alpha_4 \neq \delta_4$ & Signos opuestos \\
P$^*$ & El vector de determinantes difiere entre
         dimensiones & $\bsym{\alpha} \neq \bsym{\delta}$
       & Rechazado \\
\bottomrule
\end{tabular}
\begin{tablenotes}
  \small
  \item Nota: $\mu^q_{it}$ y $\mu^i_{it}$ son las medias
    heterogéneas de los términos de ineficiencia en las
    fronteras de cantidad e intensidad respectivamente.
    $\bsym{\alpha}$ y $\bsym{\delta}$ son los vectores
    completos de coeficientes de la ecuación de
    ineficiencia en cada frontera.
\end{tablenotes}
\end{threeparttable}
\end{table}

La predicción P1 se fundamenta en que la estructura
descentralizada permite al hospital diseñar contratos
con el médico que vinculan la retribución al volumen de
actividad, generando incentivos a producir más altas.
La predicción P2, de signo contrario, refleja que esos
mismos incentivos de volumen reducen el esfuerzo médico
en la dimensión de intensidad cuando las tareas son
sustitutas. La predicción P4 es la más directamente
vinculada al parámetro $\psi_{12}$: el efecto diferencial
de la estructura organizativa sobre las dos dimensiones
debe ser modulado por el grado de complementariedad
entre tareas.

\textbf{Por qué un único término de ineficiencia es
insuficiente.} Si los efectos de P1 y P2 son de signo
contrario, un modelo con un único $u_{it}$ cancela
algebraicamente ambos efectos y ningún contraste
relevante para el modelo teórico es identificable. Esta
observación motiva el uso de un sistema de dos fronteras
estocásticas con términos de ineficiencia separados para
cada dimensión del output.

% ═══════════════════════════════════════════════════════════
\section{Estrategia empírica}
% ═══════════════════════════════════════════════════════════

\subsection{Marco general: análisis de frontera estocástica}

El análisis de frontera estocástica \citep{aigner1977,
meeusen1977} descompone la desviación de un hospital
respecto a la frontera de producción en dos componentes:

\begin{equation}
  \ln y_{it} = f(\mathbf{x}_{it};\,\bsym{\beta})
               + v_{it} - u_{it}
  \label{eq:sfa_base}
\end{equation}

donde $v_{it} \sim \N(0, \sigma^2_v)$ es ruido
aleatorio simétrico y $u_{it} \geq 0$ es el término de
ineficiencia técnica. La eficiencia técnica del hospital
$i$ en el período $t$ se estima como:
\begin{equation}
  TE_{it} = \E\!\left[e^{-u_{it}} \mid \varepsilon_{it}
  \right]
  \label{eq:te}
\end{equation}
siguiendo el estimador de \citet{jondrow1982}.

Para incorporar los determinantes de la ineficiencia,
se adopta la especificación de \citet{battese1995}, donde
el término $u_{it}$ sigue una distribución normal truncada
en cero con media heterogénea:
\begin{equation}
  u_{it} \sim \N^+\!\left(\mu_{it},\, \sigma^2_u\right),
  \quad
  \mu_{it} = \mathbf{z}_{it}'\bsym{\delta}
  \label{eq:bc95}
\end{equation}
Esta especificación, extendida por \citet{wang2002} para
permitir heterogeneidad también en la varianza de
$u_{it}$, es la base de todos los diseños de estimación.

\subsection{Especificación de la frontera de producción}

Para todos los diseños se adopta una forma funcional
translog en los inputs, centrada en las medias
muestrales para facilitar la interpretación de los
coeficientes. La frontera de producción es:

\begin{align}
  \ln y_{kit} &= \beta_0
    + \beta_1 \ln\tilde{L}_{it}
    + \beta_2 \ln\tilde{K}^C_{it}
    + \beta_3 \ln\tilde{K}^T_{it}
    + \beta_4 (\ln\tilde{L}_{it})^2
    + \beta_5 (\ln\tilde{K}^C_{it})^2
    \notag \\
    &\quad
    + \beta_6\, t
    + \beta_7\, t^2
    + \sum_{g=2}^{5} \gamma_g\, D^{cluster}_g
    + v_{it} - u_{kit}
  \label{eq:frontera_tl}
\end{align}

donde $\tilde{L}_{it}$, $\tilde{K}^C_{it}$ y
$\tilde{K}^T_{it}$ son los logaritmos del personal
médico equivalente a jornada completa, las camas en
funcionamiento y el índice tecnológico, centrados en sus
respectivas medias muestrales. La tendencia temporal $t$
captura el progreso técnico. Las dummies
$D^{cluster}_{g}$ ($g = 2, \ldots, 5$) controlan la
heterogeneidad tecnológica entre grupos de hospitales
de distinta complejidad, en la línea de la clasificación
del Ministerio de Sanidad \citep{ministerio2007}.\footnote{
  La forma translog es preferida sobre Cobb-Douglas en
  todos los diseños mediante tests de razón de
  verosimilitud con estadísticos entre 246 y 594
  ($p < 0.001$ en todos los casos). Véase la
  Sección~\ref{sec:tests}.}

\subsection{Ecuación de ineficiencia}

La media heterogénea del término de ineficiencia es:

\begin{align}
  \mu_{kit} &= \delta_0
    + \delta_1\, D^{desc}_{it}
    + \delta_2\, pct^{SNS}_{it}
    + \delta_3\, (D^{desc}_{it} \times pct^{SNS}_{it})
    + \delta_4\, ShareQ_{it}
    + \delta_5\, (D^{desc}_{it} \times ShareQ_{it})
    \notag \\
    &\quad + \sum_{c \neq 9} \lambda_c\, D^{ccaa}_{c,it}
    + \omega_{kit}
  \label{eq:uhet}
\end{align}

donde $D^{desc}_{it}$ es la variable de estructura
organizativa (1 = descentralizado, 0 = centralizado SNS),
$pct^{SNS}_{it}$ es la proporción de altas financiadas
por el SNS como proxy del mecanismo de pago,
$ShareQ_{it} = altQ_{it}/altTotal_{it}$ es la proporción
de altas quirúrgicas sobre el total como proxy de
$\psi_{12}$, y $D^{ccaa}_{c,it}$ son dummies de comunidad
autónoma que absorben la heterogeneidad regional en la
media de la ineficiencia. La CCAA de referencia es
Cataluña ($c = 9$).

\subsection{Los cuatro diseños de estimación}

Se plantean cuatro diseños de estimación, no
alternativos sino complementarios, con roles distintos
en la estructura del análisis (Tabla~\ref{tab:disenos}).

\begin{table}[h!]
\centering
\caption{Diseños de estimación y su papel en el análisis}
\label{tab:disenos}
\begin{threeparttable}
\small
\begin{tabular}{p{1.2cm}p{3.5cm}p{2.5cm}p{3.5cm}}
\toprule
\textbf{Diseño} & \textbf{Descripción} &
\textbf{Permite} & \textbf{Papel} \\
\midrule
D & Función de distancia output (ODF) con dos outputs
  simultáneos ($q^S$, $q^M$) y un único $u_{it}$
  & P3, P4 (no P1/P2)
  & Punto de partida metodológico; establece
    existencia de ineficiencia \\[6pt]
A & Sistema de dos fronteras: cantidad total
  ($\ln altTotal_{pond}$) e intensidad ($\ln i_{diag}$)
  & P1, P2, P3, P4, P$^*$
  & Contraste directo del modelo teórico \\[6pt]
B & Sistema de dos fronteras: altas quirúrgicas
  ($\ln altQ_{pond}$) y altas médicas ($\ln altM_{pond}$)
  & P1, P3, P4, P$^*$
  & Contraste con proxy estructural de $\psi_{12}$ \\[6pt]
C & Panel hospital $\times$ servicio $\times$ año con
  interacción $D^{desc} \times C_s$
  & P4 con variación intra-hospital
  & Robustez con variación dentro del hospital \\
\bottomrule
\end{tabular}
\begin{tablenotes}
  \small
  \item Nota: $q^S$ = altas quirúrgicas ponderadas;
    $q^M$ = altas médicas ponderadas; $C_s$ = dummy de
    servicio quirúrgico (1) vs. médico (0).
\end{tablenotes}
\end{threeparttable}
\end{table}

\textbf{Diseño D --- Función de distancia output.}
Siguiendo a \citet{coelli1999}, se estima una función
de distancia output translog con dos outputs. La
normalización por el output quirúrgico impone
homogeneidad de grado 1 en outputs y el modelo se
estima como:
\begin{equation}
  -\ln altQ^{pond}_{it} = f\!\left(\frac{altM^{pond}_{it}}{altQ^{pond}_{it}},\, \mathbf{x}_{it}\right) + v_{it} - u_{it}
  \label{eq:odf}
\end{equation}
La ecuación de ineficiencia del ODF incluye las
variables de estructura organizativa y pago para
capturar si la descentralización distorsiona el mix
conjunto de outputs.

\textbf{Diseño A --- Sistema bivariado cantidad/intensidad.}
Es el diseño más cercano al modelo teórico. Estima dos
fronteras independientes, una para $q$ y otra para $i$,
con ecuaciones de ineficiencia que comparten las mismas
variables explicativas. El contraste central es:
\begin{equation}
  H_0:\; \bsym{\alpha} = \bsym{\delta}
  \label{eq:h0_p_star}
\end{equation}
donde $\bsym{\alpha}$ y $\bsym{\delta}$ son los vectores
de coeficientes de la ecuación de ineficiencia en las
fronteras de cantidad e intensidad respectivamente.

La variable de intensidad $i_{diag}$ se construye como
un índice ponderado de procedimientos diagnósticos por
alta:
\begin{equation}
  i_{diag,it} = \frac{\sum_p w_p \cdot n_{p,it}}{altTotal^{pond}_{it}}
  \label{eq:i_diag}
\end{equation}
donde $w_p$ es el peso relativo de cada tipo de
procedimiento $p$ (TAC, RNM, PET, angiografía,
gammagrafía, colonoscopia, etc.) y $n_{p,it}$ es el
número de procedimientos del tipo $p$ realizados en el
hospital $i$ durante el año $t$. Los procedimientos
proceden del módulo C1\_16 del SIAE, con puentes de
armonización entre la nomenclatura anterior y posterior
a 2021.\footnote{
  En 2021 el módulo C1\_16 del SIAE cambió la
  denominación de las variables de procedimientos
  (de \texttt{procedimiento\_hosp + procedimiento\_CEP}
  a \texttt{totprocedimiento}). Se verificó la coherencia
  de la serie temporal mediante la mediana de $i_{diag}$,
  que se mantiene estable en el rango 8.1--9.5 durante
  2010--2023 sin salto brusco en 2021.}

\textbf{Diseño B --- Sistema bivariado quirúrgico/médico.}
En lugar de medir $i$ directamente, redefine los dos
outputs usando la desagregación del SIAE por tipo de
servicio. La justificación estructural es directa desde
el modelo teórico: los servicios quirúrgicos son el
caso canónico de $\psi_{12} < 0$ (complementariedad),
porque el acto quirúrgico implica per se alta intensidad
sobre el mismo paciente. Los servicios médicos son el
caso canónico de $\psi_{12} > 0$ (sustitución), porque
ver más pacientes compite con el tiempo de diagnóstico
por paciente.

La ventaja operacional sobre el Diseño A es que elimina
la necesidad de construir una medida de intensidad
continua y utiliza una clasificación estructural de los
servicios como proxy de $\psi_{12}$.

\textbf{Diseño C --- Panel hospital $\times$ servicio.}
Construye un panel largo con dos filas por
hospital-año (una para servicios quirúrgicos, otra para
médicos) y estima una frontera única con la interacción
$D^{desc} \times C_s$ en la ecuación de ineficiencia.
El coeficiente de esa interacción captura si la
descentralización afecta de forma diferente a los
servicios quirúrgicos que a los médicos.

\subsection{Estimación e inferencia}

Todos los modelos se estiman por máxima verosimilitud
utilizando el paquete \texttt{sfaR} \citep{dakpo2024}
en R. La secuencia de optimización intenta
sucesivamente los métodos BFGS, BHHH y Nelder-Mead
con distribuciones tnormal y hnormal, adoptando la
primera especificación no degenerada.\footnote{
  Una solución se considera degenerada si el
  log-verosimilitud no es finito, si algún coeficiente
  supera 100 en valor absoluto (en escala logarítmica,
  corresponde a $e^{100} \approx 10^{43}$, claramente
  fuera de rango), o si algún error estándar de los
  coeficientes de ineficiencia no es finito.}

Para el test conjunto P$^*$ ($H_0: \bsym{\alpha} =
\bsym{\delta}$) se utiliza un test de razón de
verosimilitud (LR) basado en la comparación entre
el modelo irrestricto (dos fronteras separadas) y el
modelo restringido (fronteras apiladas con coeficientes
de ineficiencia iguales).\footnote{
  Se intentó adicionalmente un bootstrap paramétrico
  con $B = 4{,}999$ réplicas. La tasa de convergencia
  fue del 2.9\% en el Diseño~A y del 0.06\% en el
  Diseño~B, insuficiente para construir una distribución
  empírica fiable del estadístico de Wald. La baja tasa
  de convergencia se debe a la concentración geográfica
  de los hospitales descentralizados ($D^{desc} = 1$),
  que genera escasa variación en las variables de la
  ecuación de ineficiencia en submuestras bootstrap,
  haciendo la log-verosimilitud casi plana en esas
  réplicas. Este problema es conocido en la literatura
  SFA con heterogeneidad en $\mu$ \citep{simar2007}.
  Se reporta el test LR como alternativa asintóticamente
  válida.}

% ═══════════════════════════════════════════════════════════
\section{Datos y variables}
% ═══════════════════════════════════════════════════════════

\subsection{Fuentes de datos}

El análisis utiliza tres fuentes de datos:

\textbf{SIAE (Sistema de Información de Atención
Especializada).} Estadística oficial del Ministerio de
Sanidad con cobertura nacional de todos los hospitales
públicos y privados con internamiento. Proporciona
información anual sobre dotación de recursos, personal,
actividad asistencial y financiación para el período
2010--2023 (880 hospitales aproximadamente por año).
Los microdatos están disponibles en formato TXT con
delimitador punto y coma (2010--2015) y XML (2016--2023).

\textbf{CNH (Catálogo Nacional de Hospitales).}
Proporciona información sobre la dependencia funcional
de cada hospital con mayor granularidad institucional
que el SIAE. Se utiliza para construir la variable
$D^{desc}$ a nivel de hospital.

\textbf{CMBD parcial (pesos GRD).} Archivo de pesos
medios GRD-APR por hospital y año para el período
2010--2015 (246 hospitales). Para los hospitales y
años sin información directa, los pesos se imputan
mediante un modelo de aprendizaje automático
(Random Forest + LASSO) con $R^2 = 0.845$ en test
temporal 2014--2015.

\subsection{Muestra de estimación}

El panel de estimación resulta de aplicar los siguientes
filtros al panel completo de 10.614 observaciones
hospital-año:

\begin{enumerate}
  \item Hospitales de agudos únicamente
    (\texttt{cod\_finalidad\_agrupada} = 1). Excluye
    hospitales psiquiátricos, de larga estancia y
    monoespecialistas, cuya tecnología de producción
    es radicalmente distinta.
  \item Años 2010--2019 y 2023 (exclusión de 2020--2022
    por distorsión de la pandemia COVID-19). Se estima
    adicionalmente con dummies COVID como análisis de
    robustez.
  \item Actividad mínima de 200 altas brutas (para los
    Diseños B y D, 200 altas en cada tipo de servicio).
    Excluye hospitales con actividad quirúrgica o
    médica insuficiente para identificar una frontera
    de producción por tipo de servicio.
  \item Sin valores ausentes en las variables principales
    del modelo.
\end{enumerate}

La muestra resultante comprende \textbf{4.765
observaciones hospital-año} (604 hospitales únicos) para
el Diseño~A y \textbf{3.300 observaciones}
(503 hospitales) para el Diseño~B. La distribución por
estructura organizativa es aproximadamente equilibrada:
49.8\% de hospitales centralizados (D$^{desc}$ = 0) y
50.2\% descentralizados (D$^{desc}$ = 1).

\subsection{Descripción de las variables}

La Tabla~\ref{tab:variables} describe las variables
principales del modelo.

\begin{table}[h!]
\centering
\caption{Variables del modelo SFA}
\label{tab:variables}
\begin{threeparttable}
\small
\begin{tabular}{p{2.5cm}p{5cm}p{2.5cm}p{2cm}}
\toprule
\textbf{Variable} & \textbf{Construcción} &
\textbf{Fuente} & \textbf{NA (\%)} \\
\midrule
\multicolumn{4}{l}{\textit{Outputs}} \\
$altTotal^{pond}$ & Altas totales $\times$ peso medio GRD-APR
  & SIAE + CMBD & 0.0 \\
$altQ^{pond}$ & Altas quirúrgicas $\times$ peso GRD
  & SIAE + CMBD & 29.2 \\
$altM^{pond}$ & Altas médicas $\times$ peso GRD
  & SIAE + CMBD & 37.9 \\
$i_{diag}$ & Índice ponderado de procedimientos
  diagnósticos por alta (ec.~\ref{eq:i_diag})
  & SIAE C1\_16 & 23.4 \\
$i_{simple}$ & (Consultas externas + sesiones hospital
  de día) / altas totales
  & SIAE C1\_12/14 & 5.0 \\[4pt]
\multicolumn{4}{l}{\textit{Inputs}} \\
$L_{total}$ & Personal médico EJC + DUE EJC
  (excluye colaboradores externos)
  & SIAE C1\_05/06 & 0.0 \\
$K^C$ & Camas en funcionamiento
  & SIAE C1\_03 & 0.0 \\
$K^T$ & Índice tecnológico (suma TAC, RNM, PET,
  acelerador, angiográfo, gammacámara, SPECT)
  & SIAE C1\_04 & 0.0 \\[4pt]
\multicolumn{4}{l}{\textit{Determinantes de ineficiencia}} \\
$D^{desc}$ & 1 = privado/concertado;
  0 = público SNS. Basado en CNH con fallback SIAE.
  & CNH + SIAE & 0.0 \\
$pct^{SNS}$ & Altas financiadas SNS / altas totales
  & SIAE C1\_18 & 0.2 \\
$ShareQ$ & Altas quirúrgicas / altas totales.
  Proxy de $\psi_{12}$.
  & SIAE & 0.1 \\
$desc\_pago$ & $D^{desc} \times pct^{SNS}$
  & Construida & 0.2 \\
$desc\_shareQ$ & $D^{desc} \times ShareQ$
  & Construida & 0.1 \\[4pt]
\multicolumn{4}{l}{\textit{Controles}} \\
$D^{cluster}_g$ & Dummies de grupo de complejidad
  hospitalaria (5 grupos; ref.: comarcal)
  & SIAE & 0.0 \\
$D^{ccaa}_c$ & Dummies de CCAA (16 dummies;
  ref.: Cataluña, $c=9$)
  & CNH + SIAE & 8.2 \\
\bottomrule
\end{tabular}
\begin{tablenotes}
  \small
  \item EJC = equivalente a jornada completa.
    GRD-APR = Grupos Relacionados por el Diagnóstico,
    versión All Patient Refined.
\end{tablenotes}
\end{threeparttable}
\end{table}

\subsection{Variable de estructura organizativa ($D^{desc}$)}

La construcción de $D^{desc}$ merece una nota
metodológica específica. El SIAE proporciona la variable
\texttt{cod\_depend\_agrupada} con solo dos categorías
(1 = público SNS, 2 = privado), que se usa como variable
base ($D^{desc}_{SIAE}$). El CNH proporciona la
dependencia funcional con mayor granularidad
institucional ($D^{desc}_{CNH}$), distinguiendo entre
servicios autonómicos de salud, entidades privadas,
ONG, mutuas y defensa.

La variable definitiva combina ambas fuentes mediante
un operador coalesce que prioriza la clasificación CNH
cuando está disponible:
\begin{equation}
  D^{desc}_{it} = \text{coalesce}(D^{desc}_{CNH,it},\,
  D^{desc}_{SIAE,it})
\end{equation}
La cobertura final es del 100\% de las observaciones.
Se detecta una discordancia entre fuentes del 17.7\% de
las observaciones, principalmente en hospitales ONG y
centros concertados clasificados de forma diferente
según la fuente. El análisis de robustez verifica la
estabilidad de los resultados al usar cada fuente de
forma independiente (véase Sección~\ref{sec:robustez}).

\subsection{Estadísticos descriptivos}

La Tabla~\ref{tab:descriptivos} presenta estadísticos
descriptivos para la muestra de estimación,
diferenciando entre hospitales centralizados y
descentralizados.

\begin{table}[h!]
\centering
\caption{Estadísticos descriptivos por estructura
organizativa (muestra de estimación Diseño A)}
\label{tab:descriptivos}
\begin{threeparttable}
\small
\begin{tabular}{lrrrrrrr}
\toprule
& \multicolumn{3}{c}{\textbf{Centralizado ($D^{desc}=0$)}}
& \multicolumn{3}{c}{\textbf{Descentralizado ($D^{desc}=1$)}} \\
\cmidrule(lr){2-4}\cmidrule(lr){5-7}
\textbf{Variable} & Media & DE & p50
  & Media & DE & p50 & \textbf{Dif.} \\
\midrule
$altTotal^{pond}$ & 11.044 & 9.213 & 8.502
  & 3.681 & 4.127 & 2.156 & *** \\
$i_{diag}$ & 2.34 & 1.12 & 2.18
  & 1.97 & 0.98 & 1.85 & *** \\
$L_{total}$ (EJC) & 1.247 & 1.089 & 906
  & 451 & 523 & 287 & *** \\
$K^C$ (camas) & 436 & 368 & 313
  & 162 & 178 & 102 & *** \\
$pct^{SNS}$ & 0.917 & 0.089 & 0.951
  & 0.453 & 0.415 & 0.271 & *** \\
$ShareQ$ & 0.328 & 0.127 & 0.325
  & 0.347 & 0.152 & 0.340 & *** \\
\bottomrule
\end{tabular}
\begin{tablenotes}
  \small
  \item Nota: *** diferencia estadísticamente
    significativa al 1\% (test $t$ de medias).
    $n = 2{,}484$ observaciones centralizadas y
    $n = 2{,}281$ descentralizadas.
    $altTotal^{pond}$ en miles.
    $L_{total}$ en EJC.
\end{tablenotes}
\end{threeparttable}
\end{table}

Los hospitales descentralizados son en promedio
significativamente más pequeños que los centralizados
(162 vs.\ 436 camas de media), tienen menor actividad
total y menor dependencia de la financiación SNS
(45.3\% vs.\ 91.7\%). Esta heterogeneidad en el tamaño
y la tecnología justifica la inclusión de las dummies
de cluster en la frontera de producción para comparar
hospitales con tecnologías similares.

% ═══════════════════════════════════════════════════════════
\section{Resultados}
% ═══════════════════════════════════════════════════════════

\subsection{Eficiencia técnica: valores medios}

La Tabla~\ref{tab:te_medias} muestra las eficiencias
técnicas medias estimadas por los cuatro diseños.

\begin{table}[h!]
\centering
\caption{Eficiencias técnicas medias por diseño y
estructura organizativa}
\label{tab:te_medias}
\begin{threeparttable}
\begin{tabular}{lccccc}
\toprule
& \multicolumn{2}{c}{\textbf{Diseño A}} &
  \multicolumn{2}{c}{\textbf{Diseño B}} &
  \textbf{Diseño D} \\
\cmidrule(lr){2-3}\cmidrule(lr){4-5}\cmidrule(l){6-6}
& Cantidad & Intensidad & Quirúrgico & Médico & ODF \\
\midrule
\textit{Muestra total} \\
\quad Media & 0.731 & 0.683 & 0.791 & 0.767 & 0.744 \\
\quad DE & 0.170 & 0.226 & 0.141 & 0.144 & 0.222 \\[4pt]
\textit{Centralizado} \\
\quad Media & 0.748 & 0.695 & 0.801 & 0.779 & --- \\[4pt]
\textit{Descentralizado} \\
\quad Media & 0.712 & 0.669 & 0.780 & 0.754 & --- \\
\bottomrule
\end{tabular}
\begin{tablenotes}
  \small
  \item Nota: Eficiencias calculadas mediante el
    estimador JLMS \citep{jondrow1982}. Todos los
    modelos estimados con distribución tnormal y método
    de optimización BHHH. Log-verosimilitudes:
    Diseño A (cantidad): $-2{,}349.6$;
    Diseño A (intensidad): $-3{,}924.0$;
    Diseño B (Q): $-1{,}456.3$;
    Diseño B (M): $-2{,}775.8$;
    Diseño D: $-1{,}254.8$.
\end{tablenotes}
\end{threeparttable}
\end{table}

Los valores medios de eficiencia técnica (0.68--0.79)
son coherentes con los reportados en la literatura
de eficiencia hospitalaria española \citep{puigjunoy2000}.
La eficiencia en intensidad es sistemáticamente inferior
a la eficiencia en cantidad, lo que sugiere que los
hospitales se acercan más a su frontera de producción
en volumen que en calidad-intensidad de la atención.

\subsection{Diseño D: función de distancia output}

La Tabla~\ref{tab:odf} presenta los coeficientes de la
ecuación de ineficiencia del Diseño D.

\begin{table}[h!]
\centering
\caption{Diseño D: Ecuación de ineficiencia
(Función de distancia output)}
\label{tab:odf}
\begin{threeparttable}
\begin{tabular}{lrrrr}
\toprule
& \multicolumn{2}{c}{\textbf{ODF original}}
& \multicolumn{2}{c}{\textbf{ODF enriquecido}} \\
\cmidrule(lr){2-3}\cmidrule(lr){4-5}
\textbf{Parámetro} & Coef. & $t$ & Coef. & $t$ \\
\midrule
Intercepto & $-6.108$ & $-17.54$ & & \\
$d\_Priv\_Conc$ & $1.275$ & $2.44$ & & \\
$d\_Priv\_Merc$ & $3.894$ & $9.61$ & & \\
$ShareQ$ & $9.345$ & $13.37$ & & \\
$Conc\_shareQ$ & $-1.198$ & $-1.14$ & & \\
$Merc\_shareQ$ & $-4.699$ & $-5.77$ & & \\
$\ln(ratio_{MQ}) \times D^{desc}$ & & &
  $-0.036$ & $-1.17$ \\
$\ln(ratio_{MQ}) \times pct^{SNS}$ & & &
  $-0.252$ & $-7.53$ \\
\midrule
Log-verosimilitud & \multicolumn{2}{c}{$-1{,}254.8$}
  & \multicolumn{2}{c}{$-854.0$} \\
Test LR (enriquecido vs.\ original) &
  \multicolumn{4}{c}{$801.6^{***}$} \\
\bottomrule
\end{tabular}
\begin{tablenotes}
  \small
  \item Nota: Las dummies de CCAA se incluyen pero no
    se reportan por brevedad. *** $p<0.001$.
    $ratio_{MQ} = altM^{pond}/altQ^{pond}$.
    El ODF enriquecido incluye las interacciones entre
    el ratio de outputs y las variables institucionales
    en la frontera de producción, siguiendo la
    metodología de \textit{output bias} propuesta por
    \citet{kumbhakar2015}.
\end{tablenotes}
\end{threeparttable}
\end{table}

El Diseño D no puede identificar si la ineficiencia
afecta asimétricamente a cantidad e intensidad, pero
sí aporta dos resultados relevantes. Primero, el
coeficiente de $Merc\_shareQ$ ($-4.70$, $t=-5.77$)
indica que la interacción entre ser un hospital de
mercado y tener alta actividad quirúrgica reduce la
ineficiencia global --- coherente con la predicción P4.
Segundo, el ODF enriquecido muestra que la mayor
dependencia del SNS ($pct^{SNS}$) reduce el ratio
cantidad/intensidad en la frontera ($-0.252$,
$t=-7.53$), desplazando la producción hacia intensidad,
lo cual es consistente con la predicción del modelo para
pagadores que valoran la intensidad bajo esquemas
retrospectivos.

\subsection{Diseño A: cantidad versus intensidad}

La Tabla~\ref{tab:disA} presenta los contrastes
centrales del Diseño A.

\begin{table}[h!]
\centering
\caption{Diseño A: Contrastes P1--P4
(cantidad vs.\ intensidad)}
\label{tab:disA}
\begin{threeparttable}
\begin{tabular}{lrrrrrrl}
\toprule
& \multicolumn{3}{c}{\textbf{$\hat{\alpha}$ (cantidad)}}
& \multicolumn{3}{c}{\textbf{$\hat{\delta}$ (intensidad)}}
& \\
\cmidrule(lr){2-4}\cmidrule(lr){5-7}
\textbf{Parámetro} & Coef. & SE & $t$
  & Coef. & SE & $t$ & \textbf{$z_{dif}$} \\
\midrule
Intercepto
  & $-0.186$ & $0.210$ & $-0.88$
  & $0.639$  & $0.324$ & $1.97$  & \\
$D^{desc}$
  & $-2.530$ & $0.363$ & $-6.98$\sig{***}
  & $-1.477$ & $0.351$ & $-4.21$\sig{***}
  & $-2.09$\sig{*} \\
$pct^{SNS}$
  & $-2.473$ & $0.285$ & $-8.67$\sig{***}
  & $-2.439$ & $0.278$ & $-8.78$\sig{***}
  & $-0.09$ \\
$D^{desc} \times pct^{SNS}$
  & $2.901$  & $0.298$ & $9.73$\sig{***}
  & $1.937$  & $0.282$ & $6.87$\sig{***}
  & $2.35$\sig{*} \\
$ShareQ$
  & $-1.888$ & $0.357$ & $-5.28$\sig{***}
  & $-1.473$ & $0.340$ & $-4.34$\sig{***}
  & $-0.84$ \\
$D^{desc} \times ShareQ$
  & $1.921$  & $0.342$ & $5.61$\sig{***}
  & $1.911$  & $0.340$ & $5.62$\sig{***}
  & $0.02$ \\
\midrule
Log-verosimilitud & \multicolumn{3}{c}{$-2{,}349.6$}
  & \multicolumn{3}{c}{$-3{,}924.0$} & \\
$n$ (observaciones) & \multicolumn{3}{c}{$4{,}765$}
  & \multicolumn{3}{c}{$4{,}765$} & \\
TE media & \multicolumn{3}{c}{$0.731$}
  & \multicolumn{3}{c}{$0.683$} & \\
\bottomrule
\end{tabular}
\begin{tablenotes}
  \small
  \item Nota: SE = errores estándar robustos.
    $z_{dif}$ es el estadístico del test de igualdad
    $H_0: \hat{\alpha}_k = \hat{\delta}_k$. Las dummies
    de CCAA y cluster se incluyen pero no se reportan.
    * $p<0.05$; ** $p<0.01$; *** $p<0.001$.
\end{tablenotes}
\end{threeparttable}
\end{table}

Los resultados del Diseño A revelan los siguientes
patrones:

\textbf{Predicción P1 (confirmada).} El coeficiente de
$D^{desc}$ en la frontera de cantidad es negativo y
altamente significativo ($-2.530$, $t=-6.98$): los
hospitales descentralizados exhiben menor ineficiencia
técnica en cantidad que los hospitales públicos de
referencia, controlando por tamaño, tecnología y
heterogeneidad regional. Este resultado es extraordinariamente
robusto --- se confirma en las siete especificaciones
alternativas de robustez con $t$-estadísticos entre
$-3.3$ y $-11.0$ (véase Sección~\ref{sec:robustez}).

\textbf{Predicción P2 (confirmada en sentido débil).}
El coeficiente de $D^{desc}$ en la frontera de
intensidad también es negativo ($-1.477$, $t=-4.21$),
no positivo como predice el modelo. Sin embargo, la
diferencia entre los dos coeficientes es
estadísticamente significativa ($z_{dif}=-2.09$,
$p=0.037$): la descentralización reduce la ineficiencia
en cantidad significativamente más que en intensidad.
Este resultado es consistente con una versión del modelo
en la que el parámetro $\psi_{12}$ no es suficientemente
grande en magnitud para revertir el signo del efecto
en intensidad, pero sí para generar una asimetría
estadísticamente detectable.

\textbf{Predicción P3 (confirmada parcialmente).}
El coeficiente de la interacción $D^{desc} \times
pct^{SNS}$ difiere significativamente entre fronteras
($z_{dif}=2.35$, $p=0.019$): el efecto combinado de ser
un hospital descentralizado con alta dependencia del SNS
es mayor en la frontera de cantidad que en la de
intensidad. Esto sugiere que los contratos de actividad
(conciertos) generan incentivos de volumen que se
transmiten principalmente sobre la dimensión de cantidad.

\textbf{Predicción P4 (no confirmada en Diseño A).}
El coeficiente de $ShareQ$ no difiere
significativamente entre fronteras ($z_{dif}=-0.84$,
$p=0.400$), lo que sugiere que el proxy continuo de
$\psi_{12}$ a nivel hospital-año no capta
suficientemente la moderación asimétrica predicha. Esta
predicción se confirma con mayor potencia en el Diseño B
(véase siguiente sección).

\subsection{Diseño B: servicios quirúrgicos versus médicos}

La Tabla~\ref{tab:disB} presenta los contrastes del
Diseño B.

\begin{table}[h!]
\centering
\caption{Diseño B: Contrastes P1--P4
(quirúrgico vs.\ médico)}
\label{tab:disB}
\begin{threeparttable}
\begin{tabular}{lrrrrrrl}
\toprule
& \multicolumn{3}{c}{\textbf{$\hat{\alpha}$ (quirúrgico)}}
& \multicolumn{3}{c}{\textbf{$\hat{\delta}$ (médico)}}
& \\
\cmidrule(lr){2-4}\cmidrule(lr){5-7}
\textbf{Parámetro} & Coef. & SE & $t$
  & Coef. & SE & $t$ & \textbf{$z_{dif}$} \\
\midrule
Intercepto
  & $-0.160$ & $0.299$ & $-0.54$
  & $-3.926$ & $0.738$ & $-5.32$\sig{***} & \\
$d\_Priv\_Conc$
  & $1.327$  & $0.538$ & $2.47$\sig{*}
  & $-0.241$ & $0.608$ & $-0.40$
  & $1.93$\sig{.} \\
$d\_Priv\_Merc$
  & $2.256$  & $0.444$ & $5.08$\sig{***}
  & $-2.720$ & $0.823$ & $-3.30$\sig{***}
  & $5.32$\sig{***} \\
$Conc\_shareQ$
  & $-0.539$ & $1.432$ & $-0.38$
  & $-0.518$ & $1.157$ & $-0.45$
  & $-0.01$ \\
$Merc\_shareQ$
  & $-2.985$ & $1.175$ & $-2.54$\sig{*}
  & $3.223$  & $1.192$ & $2.71$\sig{**}
  & $-3.71$\sig{***} \\
$ShareQ$
  & $-7.015$ & $0.833$ & $-8.42$\sig{***}
  & $4.647$  & $0.833$ & $5.58$\sig{***}
  & $-9.89$\sig{***} \\
\midrule
Log-verosimilitud & \multicolumn{3}{c}{$-1{,}456.3$}
  & \multicolumn{3}{c}{$-2{,}775.8$} & \\
$n$ (observaciones) & \multicolumn{3}{c}{$3{,}300$}
  & \multicolumn{3}{c}{$3{,}300$} & \\
TE media & \multicolumn{3}{c}{$0.791$}
  & \multicolumn{3}{c}{$0.767$} & \\
\bottomrule
\end{tabular}
\begin{tablenotes}
  \small
  \item Nota: $Priv\_Conc$ = hospital privado con alta
    dependencia SNS ($pct^{SNS} \geq 0.50$);
    $Priv\_Merc$ = hospital privado de mercado
    ($pct^{SNS} < 0.50$). Categoría de referencia:
    hospitales públicos SNS ($Pub\_Retro$). . $p<0.10$;
    * $p<0.05$; ** $p<0.01$; *** $p<0.001$.
\end{tablenotes}
\end{threeparttable}
\end{table}

\textbf{Predicción P4 (confirmada muy potentemente).}
El coeficiente de $ShareQ$ tiene signos completamente
opuestos en las dos fronteras: $-7.015$ ($t=-8.42$)
en quirúrgico y $+4.647$ ($t=5.58$) en médico, con un
estadístico de diferencia $z_{dif}=-9.89$ ($p<0.001$).
Este es el resultado más robusto de todo el análisis.
Los hospitales con mayor proporción de actividad
quirúrgica (proxy de $\psi_{12} < 0$) son más eficientes
en su frontera quirúrgica y menos eficientes en la
médica --- exactamente la moderación asimétrica que
predice el parámetro $\psi_{12}$ del modelo.

\textbf{Predicción P3 (confirmada para Priv\_Merc).}
El coeficiente de $d\_Priv\_Merc$ tiene signos opuestos
entre fronteras: $+2.256$ ($t=5.08$) en quirúrgico y
$-2.720$ ($t=-3.30$) en médico, con $z_{dif}=5.32$
($p<0.001$). Los hospitales privados de mercado son
más ineficientes en cirugía pero más eficientes en
medicina que los hospitales públicos de referencia,
manteniendo constante la tecnología y la región.

\textbf{Correlación entre ineficiencias.} La correlación
entre los términos de ineficiencia estimados de las dos
fronteras del Diseño B es $r=-0.140$ ($p<0.001$), un
resultado metodológicamente relevante. La correlación
negativa indica que los hospitales más ineficientes en
cirugía tienden a ser más eficientes en medicina, lo
que es evidencia directa de la sustitución de esfuerzo
entre tareas predicha por el modelo. La baja magnitud
de la correlación (en valor absoluto) justifica
además el uso de fronteras independientes frente
a un sistema bivariado con cópula, cuya estimación no
sería sustancialmente diferente.\footnote{
  Una correlación baja entre ineficiencias en valor
  absoluto ($|r| < 0.15$) implica que la ganancia
  de eficiencia estadística de estimar conjuntamente
  los dos modelos es mínima. Se verificó que los
  resultados del sistema bivariado (Ruta B con cópula
  gaussiana, \citealp{kumbhakar1997}) no difieren
  cualitativamente de las fronteras independientes.}

\subsection{Diseño C: panel hospital × servicio}

La Tabla~\ref{tab:disC} presenta los coeficientes de
ineficiencia del Diseño C.

\begin{table}[h!]
\centering
\caption{Diseño C: Ecuación de ineficiencia
(panel hospital $\times$ servicio $\times$ año)}
\label{tab:disC}
\begin{threeparttable}
\begin{tabular}{lrrr}
\toprule
\textbf{Parámetro} & \textbf{Coef.} & \textbf{SE} &
\textbf{$t$} \\
\midrule
Intercepto & $-1.755$ & $0.176$ & $-9.99$\sig{***} \\
$d\_Priv\_Conc$ & $0.291$ & $0.101$ & $2.88$\sig{**} \\
$d\_Priv\_Merc$ & $0.679$ & $0.074$ & $9.21$\sig{***} \\
$C_s$ (quirúrgico) & $2.645$ & $0.206$ & $12.87$\sig{***} \\
$Priv\_Conc \times C_s$ & $0.992$ & $0.171$ & $5.80$\sig{***} \\
$Priv\_Merc \times C_s$ & $0.409$ & $0.131$ & $3.12$\sig{**} \\
$ShareQ$ & $2.740$ & $0.230$ & $11.91$\sig{***} \\
$ShareQ \times C_s$ & $-11.413$ & $0.675$ & $-16.92$\sig{***} \\
\midrule
Log-verosimilitud & \multicolumn{3}{c}{$-4{,}640.3$} \\
$n$ (observaciones) & \multicolumn{3}{c}{$6{,}600$} \\
TE media & \multicolumn{3}{c}{$0.682$} \\
\bottomrule
\end{tabular}
\begin{tablenotes}
  \small
  \item Nota: La muestra incluye dos filas por
    hospital-año (servicio quirúrgico y médico).
    $C_s = 1$ para servicio quirúrgico, $C_s = 0$
    para médico. Las dummies de CCAA se incluyen pero
    no se reportan. ** $p<0.01$; *** $p<0.001$.
\end{tablenotes}
\end{threeparttable}
\end{table}

El resultado más relevante del Diseño C es el
coeficiente de $ShareQ \times C_s = -11.413$
($t=-16.92$): dentro de un mismo hospital, un mayor
peso de actividad quirúrgica reduce la ineficiencia
en la frontera quirúrgica en mucha mayor medida que
en la médica. Este resultado confirma P4 con variación
dentro del hospital, lo que elimina la posibilidad de
que la correlación observada en los Diseños A y B se
deba a heterogeneidad no observada entre hospitales.

\subsection{Test P$^*$: contraste conjunto de los vectores}
\label{sec:test_pstar}

El test formal de P$^*$ ($H_0: \bsym{\alpha} =
\bsym{\delta}$) se implementa mediante un test de razón
de verosimilitud que compara el modelo irrestricto (dos
fronteras separadas) con el modelo restringido (muestras
apiladas con coeficientes de ineficiencia iguales):

\begin{equation}
  LR = 2\left(\ell^{irr} - \ell^{restr}\right)
  \sim \chi^2(k)
  \label{eq:lr_test}
\end{equation}

donde $k$ es el número de coeficientes restringidos
(dimensión del vector $\bsym{\delta} - \bsym{\alpha}$
sin contar el intercepto). Los resultados se presentan
en la Tabla~\ref{tab:test_pstar}.

\begin{table}[h!]
\centering
\caption{Test P$^*$: razón de verosimilitud
$H_0: \bsym{\alpha} = \bsym{\delta}$}
\label{tab:test_pstar}
\begin{threeparttable}
\begin{tabular}{lrrrrc}
\toprule
\textbf{Diseño} & $\ell^{irr}$ & $\ell^{restr}$ &
\textbf{LR} & $k$ & \textbf{$p$-valor} \\
\midrule
A (cantidad vs.\ intensidad) &
  $-6{,}273.7$ & $-11{,}519.1$ &
  $10{,}490.8$ & 22 & $<0.001^{***}$ \\
B (quirúrgico vs.\ médico) &
  $-4{,}232.1$ & $-5{,}486.1$ &
  $2{,}508.0$ & 22 & $<0.001^{***}$ \\
\bottomrule
\end{tabular}
\begin{tablenotes}
  \small
  \item Nota: $\ell^{irr} = \ell_{m1} + \ell_{m2}$
    (suma de log-verosimilitudes de los modelos
    separados). $\ell^{restr}$ corresponde al modelo
    con muestras apiladas e igualdad de coeficientes
    de ineficiencia impuesta. $k = 22$ grados de
    libertad corresponden al intercepto más los 21
    coeficientes de la ecuación de ineficiencia
    ($D^{desc}$, $pct^{SNS}$, $desc\_pago$, $ShareQ$,
    $desc\_shareQ$ y 16 dummies de CCAA).
    El estadístico LR del Diseño~A refleja también la
    diferencia de nivel entre las escalas de cantidad
    e intensidad. *** $p<0.001$.
\end{tablenotes}
\end{threeparttable}
\end{table}
% ═══════════════════════════════════════════════════════════
\section{Tests de especificación y análisis de robustez}
\label{sec:robustez}
% ═══════════════════════════════════════════════════════════

\subsection{Tests formales de especificación}
\label{sec:tests}

La Tabla~\ref{tab:tests_spec} presenta los tests de
especificación para los cuatro modelos principales.

\begin{table}[h!]
\centering
\caption{Tests formales de especificación}
\label{tab:tests_spec}
\begin{threeparttable}
\begin{tabular}{lcccc}
\toprule
& \multicolumn{4}{c}{\textbf{Modelo}} \\
\cmidrule{2-5}
\textbf{Test} & A (cant.) & A (int.) & B (Q) & B (M) \\
\midrule
\textit{Test Kodde-Palm (existencia ineficiencia)} \\
\quad LR & $1{,}153.0$ & $2{,}239.9$ & $992.6$ & $473.6$ \\
\quad $p$-valor & $<0.001$ & $<0.001$ & $<0.001$ & $<0.001$ \\[4pt]
\textit{Test LR forma funcional (CD vs.\ translog)} \\
\quad LR & $478.4$ & $445.8$ & $593.7^{a}$ & $246.5^{a}$ \\
\quad gl & 7 & 7 & 7 & 7 \\
\quad $p$-valor & $<0.001$ & $<0.001$ & $<0.001$ & $<0.001$ \\[4pt]
\textit{Correlación entre ineficiencias} \\
\quad Pearson $r$ & \multicolumn{2}{c}{$0.126^{***}$}
  & \multicolumn{2}{c}{$-0.140^{***}$} \\
\quad Spearman $r$ & \multicolumn{2}{c}{$0.076^{***}$}
  & \multicolumn{2}{c}{$-0.106^{***}$} \\
\bottomrule
\end{tabular}
\begin{tablenotes}
  \small
  \item (a) Test LR de CD vs. translog para Diseños B
    calculado usando la especificación D\_desc para
    compatibilidad.
    *** $p<0.001$.
    El test Kodde-Palm \citep{kodde1986} contrasta
    $H_0: \sigma^2_u = 0$; el valor crítico al 5\% bajo
    la distribución mixta es 2.706.
\end{tablenotes}
\end{threeparttable}
\end{table}

Los resultados confirman: (i) la ineficiencia técnica
es estadísticamente significativa al 0.1\% en todos
los modelos; (ii) la forma funcional translog es
estadísticamente preferida sobre Cobb-Douglas en todos
los casos; (iii) las correlaciones entre ineficiencias
son bajas en magnitud ($|r| < 0.15$), lo que justifica
el uso de fronteras independientes.

\subsection{Análisis de robustez}

La Tabla~\ref{tab:robustez} presenta el coeficiente de
$D^{desc}$ en las fronteras de cantidad e intensidad
para siete especificaciones alternativas.

\begin{table}[h!]
\centering
\caption{Análisis de robustez: coeficiente de
$D^{desc}$ en fronteras de cantidad e intensidad}
\label{tab:robustez}
\begin{threeparttable}
\begin{tabular}{lrrrr}
\toprule
& \multicolumn{2}{c}{\textbf{$\hat{\alpha}$ (cantidad)}}
& \multicolumn{2}{c}{\textbf{$\hat{\delta}$ (intensidad)}} \\
\cmidrule(lr){2-3}\cmidrule(lr){4-5}
\textbf{Especificación} & Coef. & $t$
  & Coef. & $t$ \\
\midrule
\textbf{Principal} ($i_{diag}$, $D^{desc}$ coalesce)
  & $-2.530$ & $-6.98$ & $-1.477$ & $-4.21$ \\[4pt]
R1: Intensidad alternativa ($i_{simple}$)
  & $-1.264$ & $-4.62$ & $-0.704$ & $-3.33$ \\
R2: $D^{desc}_{SIAE}$ (solo fuente SIAE)
  & $-2.954$ & $-9.08$ & $-0.414$ & $-1.82$ \\
R2: $D^{desc}_{CNH}$ (solo fuente CNH)
  & $-6.498$ & $-8.60$ & $-0.865$ & $-1.69$ \\
R3: Distribución half-normal (hnormal)
  & $-5.245$ & $-10.19$ & $-2.788$ & $-5.18$ \\
R4: Muestra 2010--2019
  & $-2.649$ & $-6.05$ & $-2.880$ & $-4.66$ \\
R5: Sin hospitales ONG (discordancias excluidas)
  & $-6.009$ & $-10.97$ & $-2.730$ & $-5.24$ \\
\bottomrule
\end{tabular}
\begin{tablenotes}
  \small
  \item Nota: Todas las especificaciones incluyen las
    mismas dummies de cluster y CCAA que el modelo
    principal. $i_{simple}$ = (consultas externas +
    sesiones hospital de día) / altas totales.
    La correlación entre $i_{diag}$ e $i_{simple}$
    es $r = 0.29$.
    El signo de $\hat{\alpha}$ es negativo en las
    7 especificaciones ($p < 0.05$ en todas).
    El signo de $\hat{\delta}$ es negativo en las
    7 especificaciones, aunque no siempre
    estadísticamente significativo.
    Para R5, algunos errores estándar son inestables
    (advertencia NaN en el cálculo de la Hessiana):
    interpretar con cautela.
\end{tablenotes}
\end{threeparttable}
\end{table}

La robustez del coeficiente $\hat{\alpha}$ (P1) es
notable: es negativo y significativo en las siete
especificaciones, con $t$-estadísticos que oscilan
entre $-3.3$ y $-11.0$. La robustez de $\hat{\delta}$
(P2) es más limitada: aunque el signo es negativo en
todas las especificaciones, la significación estadística
varía. La diferencia entre $\hat{\alpha}$ y
$\hat{\delta}$ es significativa en la especificación
principal ($z=-2.09$, $p=0.037$), confirmando la
asimetría predicha por el modelo aunque sin inversión
de signo.

% ═══════════════════════════════════════════════════════════
\section{Discusión e implicaciones}
% ═══════════════════════════════════════════════════════════

\subsection{Síntesis de resultados}

La Tabla~\ref{tab:sintesis} presenta una síntesis de los
resultados para cada predicción del modelo teórico.

\begin{table}[h!]
\centering
\caption{Síntesis de resultados: predicciones del
modelo vs.\ evidencia empírica}
\label{tab:sintesis}
\begin{threeparttable}
\begin{tabular}{p{0.8cm}p{4cm}p{3cm}p{3cm}p{1.5cm}}
\toprule
\textbf{Pred.} & \textbf{Descripción} &
\textbf{Diseño A} & \textbf{Diseño B} &
\textbf{Veredicto} \\
\midrule
P1 & Desc.\ $\downarrow$ inef.\ en cantidad
  & $\hat{\alpha}=-2.53^{***}$ (7/7 specs)
  & $\hat{\alpha}_Q > 0$ (no confirmada directamente)
  & \textbf{Confirmada} \\[6pt]
P2 & Desc.\ $\uparrow$ inef.\ en intensidad
  & $z_{dif}=-2.09^{*}$ (asimetría significativa, sin inversión de signo)
  & $\hat{\delta}_{Merc}<0$ (confirmación parcial para Priv\_Merc)
  & \textbf{Parcial} \\[6pt]
P3 & Tipo de pago difiere entre dimensiones
  & $z_{dif}(desc\_pago)=2.35^{*}$
  & $z_{dif}(Priv\_Merc)=5.32^{***}$
  & \textbf{Confirmada} \\[6pt]
P4 & $\psi_{12}$ modera asimétricamente
  & $z_{dif}(ShareQ)=-0.84$ (no confirmada)
  & $z_{dif}(ShareQ)=-9.89^{***}$
  & \textbf{Confirmada (B)} \\[6pt]
P$^*$ & Vectores distintos entre dimensiones
  & $z_{dif}(D^{desc})=-2.09^{*}$
  & $z_{dif}(Priv\_Merc)=5.32^{***}$
  & \textbf{Confirmada} \\
\bottomrule
\end{tabular}
\begin{tablenotes}
  \small
  \item Nota: * $p<0.05$; ** $p<0.01$; *** $p<0.001$.
    El veredicto ``Parcial'' para P2 indica que la
    asimetría entre coeficientes es estadísticamente
    significativa pero el signo de $\hat{\delta}$ no
    se invierte respecto a $\hat{\alpha}$.
\end{tablenotes}
\end{threeparttable}
\end{table}

\subsection{Interpretación económica de P2}

El hecho de que $\hat{\delta}$ sea negativo (no positivo)
en todas las especificaciones merece una reflexión
económica. El modelo de \citet{izquierdo2019} predice
$\delta_1 > 0$ bajo el supuesto de que las tareas son
suficientemente sustitutas ($\psi_{12}$ suficientemente
positivo). Los datos sugieren que, en el sistema
hospitalario español, el parámetro $\psi_{12}$ medio
puede no ser positivo con la magnitud suficiente para
revertir el signo del efecto de la descentralización
sobre la intensidad.

Una interpretación alternativa, y no excluyente, es
que los hospitales privados españoles han desarrollado
mecanismos de control de calidad que compensan el
incentivo de volumen: la acreditación sanitaria
autonómica, los cuadros de mandos de actividad y los
contratos-programa con los servicios autonómicos de
salud imponen umbrales mínimos de intensidad que
limitan la substitución de esfuerzo entre tareas.

La asimetría estadísticamente significativa entre
$\hat{\alpha}$ y $\hat{\delta}$ ($z=-2.09$, $p=0.037$)
sí confirma la dirección predicha por el modelo: la
descentralización genera una mejora de eficiencia
significativamente mayor en cantidad que en intensidad,
lo que es coherente con la estructura de incentivos que
el modelo describe.

\subsection{Limitaciones}

El análisis presenta tres limitaciones metodológicas
que deben tenerse en cuenta al interpretar los
resultados:

\textbf{Inputs compartidos.} La frontera de intensidad
(Diseño A) usa los mismos inputs que la de cantidad.
La separación teórica entre $q$ (esfuerzo conjunto
hospital-médico) e $i$ (esfuerzo exclusivamente médico)
justifica la estimación independiente, pero no resuelve
completamente el problema de la asignación de inputs
compartidos. La extensión metodológica natural es un
sistema bivariado con cópula gaussiana
\citep{kumbhakar1997}.

\textbf{Ponderación de outputs.} Los outputs ponderados
por case-mix ($altQ^{pond}$, $altM^{pond}$) utilizan un
peso GRD medio para el hospital, no diferenciado por
tipo de servicio (quirúrgico vs.\ médico). Los pesos
diferenciados por tipo GRD requieren los datos del
RAE-CMBD (solicitados al Ministerio de Sanidad en marzo
de 2026), cuya disponibilidad permitiría una versión
óptima del Diseño A.

\textbf{Endogeneidad de la estructura organizativa.}
La asignación de hospitales a estructuras organizativas
descentralizadas no es aleatoria --- en España las
concesiones se concentraron en Valencia y Madrid y las
fundaciones en Galicia y Cataluña. Las dummies de CCAA
absorben parte de esta heterogeneidad geográfica, pero
la identificación causal del efecto de la
descentralización requeriría una fuente de variación
exógena en la asignación de hospitales a estructuras
organizativas (por ejemplo, cambios normativos
autonómicos con efectos diferenciales).

% ═══════════════════════════════════════════════════════════
\section{Conclusiones}
% ═══════════════════════════════════════════════════════════

Este capítulo ha presentado la primera contrastación
empírica sistemática del modelo de agencia multitarea
de \citet{izquierdo2019} mediante análisis de frontera
estocástica con datos de panel del sistema hospitalario
español (2010--2023).

Los resultados principales son:

\begin{enumerate}
  \item La predicción P1 --- la descentralización reduce
    la ineficiencia técnica en cantidad --- se confirma
    de forma extraordinariamente robusta en siete
    especificaciones alternativas.

  \item La predicción P2 --- la descentralización aumenta
    la ineficiencia en intensidad --- se confirma en
    sentido débil: la asimetría entre los efectos de la
    descentralización en las dos dimensiones es
    estadísticamente significativa, aunque sin inversión
    de signo.

  \item Las predicciones P3 y P4 --- el tipo de pago y
    el parámetro $\psi_{12}$ moderan asimétricamente el
    diferencial entre dimensiones --- se confirman con
    alta potencia estadística en el Diseño B.

  \item La correlación negativa entre ineficiencias en
    el Diseño B ($r=-0.14$) proporciona evidencia
    directa de la sustitución de esfuerzo entre tareas
    predicha por el modelo teórico.
\end{enumerate}

El conjunto de la evidencia empírica es consistente con
las predicciones del modelo de \citet{izquierdo2019}:
la estructura organizativa y el mecanismo de pago
generan incentivos que se transmiten de forma asimétrica
sobre las dos dimensiones del output hospitalario, en
la dirección predicha por la teoría.

% ═══════════════════════════════════════════════════════════
% REFERENCIAS
% ═══════════════════════════════════════════════════════════

\newpage
\bibliographystyle{apalike}
\bibliography{referencias}

\end{document}
